% Generated draft: requires mathematical and source review.
\documentclass[11pt]{article}
\usepackage[margin=25mm]{geometry}
\usepackage{fontspec}
\setmainfont{Latin Modern Roman}
\usepackage{amsmath,amssymb,mathtools,unicode-math}
\setmathfont{Latin Modern Math}
\usepackage{enumitem,needspace,xurl,hyperref}
\setlength{\parindent}{0pt}
\setlength{\parskip}{5pt}
\setlength{\emergencystretch}{3em}
\newcommand{\N}{\mathbb{N}}
\newcommand{\Z}{\mathbb{Z}}
\newcommand{\Q}{\mathbb{Q}}
\newcommand{\R}{\mathbb{R}}
\newcommand{\C}{\mathbb{C}}
\newcommand{\F}{\mathbb{F}}
\newcommand{\A}{\mathbb{A}}
\newcommand{\PP}{\mathbb{P}}
\newcommand{\E}{\mathbb{E}}
\newcommand{\eps}{\varepsilon}
\newcommand{\dd}{\,\mathrm{d}}
\begin{document}
\section*{UnsolvedMath ID: 20002462}
\textbf{Source problem code:} catalogue entry
\subsection{Definitions and mathematical statement}
Cory Hauck's participant note describes three research directions for hydrodynamic electron-transport balance laws, with conservative terms, collisional relaxation, and electric drift coupled to Poisson's equation: distinguish discretization-induced interface-current oscillations and negative density or temperature from model defects in high-field transition regimes; devise efficient implicit methods for the stiff flux and relaxation terms in the low-field, high-density drift-diffusion limit despite CFL restrictions; and determine appropriate relaxation times and their effect on current and temperature. The note contrasts the gas-dynamic conductivity/viscosity ratio with the single mobility coefficient in electron drift-diffusion. It does not specify a single quantified conjecture.

The supplied diagnostic is explicitly restricted to one-dimensional, nondegenerate, isothermal drift-diffusion with prescribed potential and ideal interface transmission. For temperature $T>0$, Boltzmann constant $k_B$, effective density of states $N_C>0$, and energy potential $U$, define electrochemical activity
\[
 a=\frac{n}{N_C}\exp\!\left(\frac{U}{k_BT}\right).
\]
In this model zero flux is equivalent to constant activity. Across an ideal interface the reported equilibrium jump is
\[
 \frac{n_R}{n_L}=\frac{N_{C,R}}{N_{C,L}}
 \exp\!\left(-\frac{U_R-U_L}{k_BT}\right).
\]
After fixed reference units, absorb $N_C$ into a dimensionless effective potential $V$ so $a_i=n_i e^{V_i}$ and equilibria are $n_i=C e^{-V_i}$, $C>0$.

For grid spacing $h>0$, diffusivity $D>0$, and any finite jump $\Delta V=V_{i+1}-V_i$, define $B(z)=z/(e^z-1)$ with $B(0)=1$. The exponentially fitted Scharfetter--Gummel flux in this convention is
\[
 J_{i+1/2}=\frac D h\bigl[B(\Delta V)n_i-B(-\Delta V)n_{i+1}\bigr].
\]
Let $\mathcal L(x,y)=(y-x)/(\log y-\log x)$ for positive unequal $x,y$, and $\mathcal L(x,x)=x$. The logarithmic-mean shared-potential traces may be written
\[
 n_i^*=\frac{a_i}{\mathcal L(e^{V_i},e^{V_{i+1}})},\qquad
 n_{i+1}^*=\frac{a_{i+1}}{\mathcal L(e^{V_i},e^{V_{i+1}})},\qquad
 J_{i+1/2}=\frac D h(n_i^*-n_{i+1}^*).
\]
Thus for positive cell densities, zero fitted flux, equal activity, and equal positive reconstructed traces are equivalent. The traces are nonnegative for nonnegative data. At the same Boltzmann state a strict raw-density difference detector instead gives the reported residual
\[
 \frac{DC}{h}(e^{-V_i}-e^{-V_{i+1}}),
\]
which is nonzero whenever the effective potential jump is nonzero.

The aid reports that the conservative semidiscretization built from the fitted flux preserves positivity and dissipates relative entropy for the fixed-potential equilibrium, with compatible closed or periodic boundary conditions. Relative entropy here has the form $\sum_i h[n_i\log(n_i/n_i^{\mathrm{eq}})-n_i+n_i^{\mathrm{eq}}]$, with the usual continuous convention at $n_i=0$. The full boundary and coefficient hypotheses are absent from the summary and require review. This diagnostic exhibits a discretization-induced equilibrium-current mechanism; it does not classify driven high-field oscillations or solve the original stiff hydrodynamic, Poisson-coupled, or relaxation-time problems.
\subsection{Short English statement}
The source is a research agenda about semiconductor transport. Its narrow supplied result shows how an interface flux can preserve the correct equilibrium jump, while a raw-density detector produces a spurious residual. This does not decide the origin of general high-field oscillations or provide the requested hydrodynamic solver.
\subsection{Sources}
\begin{itemize}
\item Local UnsolvedMath registry, permanent ID 20002462, source code catalogue entry; created and updated 14 August 2026. Canonical \texttt{aim-pdes-notes.json}, zero-based \texttt{notes[73]}, tag \texttt{section}, item A.9, ``Hauck, Cory.''
\item AIM ``Stiff Sources and Numerical Methods for Conservation Laws,'' version dated 1 April 2005 according to the supplied verification: \url{https://aimath.org/WWN/balancelaws/balancelaws.pdf}. Original remarks and literature field are empty; wording is marked exact.
\item Research attempt 1: \texttt{context\_only}, problem status at run \texttt{not\_a\_problem}, rigor label \texttt{proved}, novelty confidence low. It names Scharfetter--Gummel fitting but supplies no complete authoritative bibliographic reference or complete semiconductor model specification.
\item Ulam AI / UnsolvedMath Contributors, CC BY 4.0. Machine-generated research aid, not peer reviewed; independently verify before citation.
\end{itemize}
\end{document}
